\paragraph{Cayley--Poisson 乘极展开与 \texorpdfstring{$\chi^2$}{chi2} 能量谱}
以下结果在任意有限均值的输入测度下给出 \(\tilde h_t/\tilde g_t\) 的精确傅里叶展开、\(\chi^2\) 与 Rényi 阶 \(2\) 的闭式能量谱，以及乘极矩序列的尺度半群与耗散恒等式。

\begin{proposition}[Cayley--Poisson 乘极展开与傅里叶矩系数]\label{prop:xi-poisson-cauchy-cayley-multipole-fourier}
\leanverified{paper\_xi\_poisson\_cauchy\_cayley\_multipole\_fourier}
设 \(\tilde\nu\in\mathrm{Prob}(\RR)\) 具有有限均值
\[
\bar\gamma:=\int_{\RR}\gamma\,\dd\tilde\nu(\gamma).
\]
对 \(t>0\) 定义 Cauchy 核与混合密度
\[
\mathsf P_t(x):=\frac{1}{\pi}\frac{t}{x^2+t^2},
\qquad
\tilde g_t(x):=\mathsf P_t(x-\bar\gamma),
\qquad
\tilde h_t(x):=(\mathsf P_t*\tilde\nu)(x)=\int_{\RR}\mathsf P_t(x-\gamma)\,\dd\tilde\nu(\gamma),
\]
并记密度比 \(u_t(x):=\tilde h_t(x)/\tilde g_t(x)\)。
定义 Cayley 相位变量
\[
\zeta_t(x):=\frac{(x-\bar\gamma)-\mathrm{i}t}{(x-\bar\gamma)+\mathrm{i}t}\in\TT,
\]
以及乘极变量
\[
W_t(\gamma):=\frac{\gamma-\bar\gamma}{(\gamma-\bar\gamma)+2\mathrm{i}t}\in\mathbb D.
\]
则对 \(\zeta=\zeta_t(x)\in\TT\) 成立严格恒等式
\[
u_t(x)
=\EE\!\left[\frac{1-|W_t(\Gamma)|^2}{|\zeta-W_t(\Gamma)|^2}\right],
\qquad \Gamma\sim\tilde\nu,
\]
并且 \(u_t\) 在 \(\TT\) 上具有全傅里叶展开
\[
u_t(\zeta)=1+\sum_{k\ge1}\Bigl(m_k(t)\,\zeta^{-k}+\overline{m_k(t)}\,\zeta^{k}\Bigr),
\qquad
m_k(t):=\EE\!\left[W_t(\Gamma)^k\right].
\]
其中 \(m_k(t)\) 对所有 \(k\ge1\) 均无需高阶矩条件即可定义。
\end{proposition}
\begin{proof}
令 \(y:=x-\bar\gamma\)、\(a:=\gamma-\bar\gamma\)。
直接化简得到
\[
\frac{\mathsf P_t(y-a)}{\mathsf P_t(y)}=\frac{y^2+t^2}{(y-a)^2+t^2}.
\]
另一方面取 \(\zeta=\frac{y-\mathrm{i}t}{y+\mathrm{i}t}\)、\(w=\frac{a}{a+2\mathrm{i}t}\)。
则
\[
1-|w|^2=\frac{4t^2}{a^2+4t^2},
\qquad
|\zeta-w|^2=\frac{4t^2\bigl((y-a)^2+t^2\bigr)}{(y^2+t^2)(a^2+4t^2)}.
\]
故
\[
\frac{1-|w|^2}{|\zeta-w|^2}=\frac{y^2+t^2}{(y-a)^2+t^2}=\frac{\mathsf P_t(y-a)}{\mathsf P_t(y)}.
\]
对 \(a=\Gamma-\bar\gamma\) 取期望即得第一式。
\par
对 \(|w|<1\)、\(|\zeta|=1\) 的 Poisson 核标准傅里叶级数为
\[
\frac{1-|w|^2}{|\zeta-w|^2}
=1+\sum_{k\ge1}\bigl(w^k\zeta^{-k}+\overline{w}^{\,k}\zeta^{k}\bigr),
\]
而 \(|W_t(\gamma)|<1\) 对所有实 \(\gamma\) 成立，故逐项取期望合法并得到所示展开。
\end{proof}

\begin{theorem}[单时刻 Poisson--Cauchy 粗粒化的 Cayley--Fourier 反演]\label{thm:xi-poisson-cauchy-cayley-fourier-inversion-single-time}
\leanverified{paper\_xi\_poisson\_cauchy\_cayley\_fourier\_inversion\_single\_time}
在命题 \ref{prop:xi-poisson-cauchy-cayley-multipole-fourier} 的设定下，固定任意 \(t>0\)。
定义单位圆上的 Cayley 坐标
\[
\xi_t(\gamma):=\frac{(\gamma-\bar\gamma)-2\mathrm{i}t}{(\gamma-\bar\gamma)+2\mathrm{i}t}\in\TT\setminus\{1\},
\qquad
\zeta_t(x):=\frac{(x-\bar\gamma)-\mathrm{i}t}{(x-\bar\gamma)+\mathrm{i}t}\in\TT,
\]
并令 \(\mu_t:=(\xi_t)_\#\tilde\nu\in\mathrm{Prob}(\TT)\)。
则由单一时刻的密度比 \(u_t(x)=\tilde h_t(x)/\tilde g_t(x)\) 可构造性地恢复 \(\tilde\nu\)，且恢复链为
\[
u_t\ \Longrightarrow\ \{m_k(t)\}_{k\ge 0}\ \Longrightarrow\ \{\widehat\mu_t(k)\}_{k\ge 0}\ \Longrightarrow\ \mu_t\ \Longrightarrow\ \tilde\nu,
\qquad m_0(t):=1,
\]
其中：
\begin{enumerate}
  \item \textbf{（Fourier 系数读出）}将 \(u_t\) 视作 \(\TT\) 上函数 \(u_t(\zeta):=u_t\bigl(\zeta_t^{-1}(\zeta)\bigr)\)，其 Fourier 系数满足
  \[
  u_t(\zeta)=1+\sum_{k\ge1}\Bigl(m_k(t)\,\zeta^{-k}+\overline{m_k(t)}\,\zeta^{k}\Bigr),
  \qquad
  m_k(t)=\int_{\RR}W_t(\gamma)^k\,\dd\tilde\nu(\gamma),
  \]
  其中 \(W_t(\gamma)=\frac{\gamma-\bar\gamma}{(\gamma-\bar\gamma)+2\mathrm{i}t}\in\mathbb D\)。
  \item \textbf{（二项反演）}记 \(\widehat\mu_t(k):=\int_{\TT}\xi^k\,\dd\mu_t(\xi)\)。
  则对任意 \(k\ge 0\) 有上三角二项变换
  \[
  m_k(t)=2^{-k}\sum_{j=0}^{k}\binom{k}{j}\widehat\mu_t(j),
  \]
  从而可由二项反演显式恢复
  \[
  \widehat\mu_t(k)=\sum_{j=0}^{k}(-1)^{k-j}\binom{k}{j}2^{j}\,m_j(t)\qquad(k\ge 0).
  \]
  \item \textbf{（推前测度反演）}映射 \(\xi_t:\RR\to\TT\setminus\{1\}\) 为双射，其逆映射为
  \[
  \Gamma_t(\xi):=\bar\gamma+2\mathrm{i}t\,\frac{1+\xi}{1-\xi}\in\RR\qquad(\xi\in\TT\setminus\{1\}),
  \]
  并且
  \[
  \tilde\nu=(\Gamma_t)_\#\mu_t.
  \]
\end{enumerate}
\end{theorem}
\begin{proof}
第 1 条已由命题 \ref{prop:xi-poisson-cauchy-cayley-multipole-fourier} 给出，其中 \(m_k(t)\) 即 \(u_t-1\) 的 Fourier 系数。
\par
为证第 2 条，注意
\[
\xi_t(\gamma)=\frac{(\gamma-\bar\gamma)-2\mathrm{i}t}{(\gamma-\bar\gamma)+2\mathrm{i}t}
\quad\Longrightarrow\quad
\frac{1+\xi_t(\gamma)}{2}=\frac{\gamma-\bar\gamma}{(\gamma-\bar\gamma)+2\mathrm{i}t}=W_t(\gamma).
\]
故对任意 \(k\ge0\)，
\[
W_t(\gamma)^k=2^{-k}(1+\xi_t(\gamma))^k
=2^{-k}\sum_{j=0}^{k}\binom{k}{j}\xi_t(\gamma)^j.
\]
对 \(\tilde\nu\) 积分并用 \(\mu_t=(\xi_t)_\#\tilde\nu\) 得到
\(
m_k(t)=2^{-k}\sum_{j=0}^{k}\binom{k}{j}\widehat\mu_t(j)
\)；
再施行二项反演即得所示显式公式。
\par
最后，由 \(\xi=\xi_t(\gamma)\) 直接解出
\(
\gamma=\bar\gamma+2\mathrm{i}t\frac{1+\xi}{1-\xi}
\)，从而 \(\Gamma_t\) 为 \(\xi_t\) 的逆映射。
因此 \(\mu_t=(\xi_t)_\#\tilde\nu\) 等价于 \(\tilde\nu=(\Gamma_t)_\#\mu_t\)。
又由于三角多项式在 \(C(\TT)\) 中稠密，\(\{\widehat\mu_t(k)\}_{k\in\ZZ}\) 唯一决定 \(\mu_t\)，而 \(\widehat\mu_t(-k)=\overline{\widehat\mu_t(k)}\) 使得仅需 \(k\ge0\) 亦可恢复全部整数矩。
故由 \(u_t\) 的单时刻数据即可唯一确定 \(\tilde\nu\)，并给出构造性反演链。证毕。
\end{proof}

\begin{proposition}[有限原子性的 Toeplitz 秩判别与 Prony 型恢复]\label{prop:xi-poisson-cauchy-toeplitz-prony-atomic-recovery}
\leanverified{paper\_xi\_poisson\_cauchy\_toeplitz\_prony\_atomic\_recovery}
沿用定理 \ref{thm:xi-poisson-cauchy-cayley-fourier-inversion-single-time} 的记号，并设 \(\{\widehat\mu_t(k)\}_{k\ge0}\) 已由单时刻 \(u_t\) 显式恢复。
对 \(n\ge0\) 定义 Hermitian Toeplitz 矩阵族
\[
\mathbf T_n:=\bigl[\widehat\mu_t(j-k)\bigr]_{j,k=0}^{n},
\qquad
\widehat\mu_t(-k):=\overline{\widehat\mu_t(k)}.
\]
则以下命题等价：
\begin{enumerate}
  \item \(\mu_t\) 为有限原子概率测度（等价地 \(\tilde\nu\) 为有限原子概率测度）；
  \item 存在 \(N\in\NN\) 使得当 \(n\) 充分大时 \(\rank(\mathbf T_n)=N\)；
  \item 存在次数 \(N\) 的首一多项式 \(Q(z)=z^N+q_{N-1}z^{N-1}+\cdots+q_0\)，使得对所有 \(k\ge0\) 成立递推关系
  \[
  \sum_{j=0}^{N}q_j\,\widehat\mu_t(k+j)=0,\qquad q_N=1.
  \]
\end{enumerate}
在上述条件成立时，\(\mu_t\) 的原子 \(\{\xi_r\}_{r=1}^{N}\subset\TT\) 恰为 \(Q\) 的根（按重数计），其权重 \(\{p_r\}_{r=1}^{N}\) 可由 Vandermonde 系统从 \(\widehat\mu_t(0),\dots,\widehat\mu_t(N-1)\) 唯一恢复，并由 \(\gamma_r=\Gamma_t(\xi_r)\) 得
\(
\tilde\nu=\sum_{r=1}^{N}p_r\,\delta_{\gamma_r}.
\)
\end{proposition}
\begin{proof}
若 \(\mu_t=\sum_{r=1}^{N}p_r\delta_{\xi_r}\)，则 \(\widehat\mu_t(k)=\sum_{r=1}^{N}p_r\xi_r^{k}\) 为有限指数和序列。
令 \(Q(z)=\prod_{r=1}^{N}(z-\xi_r)=\sum_{j=0}^{N}q_jz^j\)，则对每个 \(k\ge0\) 有
\(
\sum_{j=0}^{N}q_j\widehat\mu_t(k+j)=\sum_{r}p_r\xi_r^{k}Q(\xi_r)=0
\)，得到 (1)\(\Rightarrow\)(3)。
\par
由 (3) 可知 \(\{\widehat\mu_t(k)\}_{k\ge0}\) 落在由首 \(N\) 项张成的有限维线性空间中，故 Toeplitz（亦等价 Hankel）矩阵的列空间维数最终不超过 \(N\)，得到 (3)\(\Rightarrow\)(2)。
\par
若 (2) 成立，则存在最小线性递推多项式 \(Q\) 消去矩序列；再由 \(\mathbf T_n\succeq0\)（来自 \(\mu_t\) 的正定性）与 Carath\'eodory--Fej\'er/Prony 型论证，\(Q\) 的根落在 \(\TT\) 上且相应表示系数为正，从而 \(\mu_t\) 为有限原子测度，得到 (2)\(\Rightarrow\)(1)。
最后 \(\tilde\nu=(\Gamma_t)_\#\mu_t\) 给出原子位置的显式恢复。证毕。
\end{proof}

\begin{proposition}[\texorpdfstring{$\chi^2$}{chi2} 散度的谱恒等式]\label{prop:xi-poisson-cauchy-chi2-spectrum-multipole}
\leanverified{paper\_xi\_poisson\_cauchy\_chi2\_spectrum\_multipole}
在命题 \ref{prop:xi-poisson-cauchy-cayley-multipole-fourier} 的设定下，令
\[
\chi^2(\tilde h_t\mid\tilde g_t)
:=\int_{\RR}\frac{(\tilde h_t-\tilde g_t)^2}{\tilde g_t}\,\dd x
=\int_{\RR}\tilde g_t(x)\,(u_t(x)-1)^2\,\dd x.
\]
则有严格恒等式
\[
\boxed{\;
\chi^2(\tilde h_t\mid\tilde g_t)=2\sum_{k\ge1}|m_k(t)|^2\; }.
\]
\end{proposition}
\begin{proof}
由 Cayley 变量 \(\zeta_t\) 的 Haar 回拉恒等式（命题 \ref{prop:app-busemann-poisson-minus1} 的等价表述），测度 \(\tilde g_t(x)\,\dd x\) 推前为 \(\TT\) 上 Haar 测度 \(\dd\theta/(2\pi)\)。
因此
\[
\chi^2(\tilde h_t\mid\tilde g_t)=\int_{0}^{2\pi}\bigl(u_t(e^{\mathrm{i}\theta})-1\bigr)^2\,\frac{\dd\theta}{2\pi}.
\]
由命题 \ref{prop:xi-poisson-cauchy-cayley-multipole-fourier}，\(u_t-1\) 的傅里叶系数满足 \(c_{-k}=m_k(t)\)、\(c_k=\overline{m_k(t)}\)（\(k\ge1\)），而 \(c_0=0\)。
Parseval 恒等式给出
\[
\int_{0}^{2\pi}\bigl(u_t-1\bigr)^2\,\frac{\dd\theta}{2\pi}
=\sum_{k\in\ZZ}|c_k|^2
=2\sum_{k\ge1}|m_k(t)|^2.
\]
\end{proof}

\begin{corollary}[Rényi 阶 \(2\) 的闭式表示与对 KL 的非渐近上界]\label{cor:xi-poisson-cauchy-renyi2-kl-upper}\leanverified{paper\_xi\_poisson\_cauchy\_renyi2\_kl\_upper}
在上述设定下，Rényi 阶 \(2\) 相对熵满足闭式
\[
D_2(\tilde h_t\mid\tilde g_t)=\log\!\left(1+2\sum_{k\ge1}|m_k(t)|^2\right).
\]
并且由 \(D_{\mathrm{KL}}\le D_2\) 得到非渐近上界
\[
D_{\mathrm{KL}}(\tilde h_t\mid\tilde g_t)\le \log\!\left(1+2\sum_{k\ge1}|m_k(t)|^2\right).
\]
\end{corollary}
\begin{proof}
由定义
\[
D_2(\tilde h_t\mid\tilde g_t)=\log\int_{\RR}\frac{\tilde h_t(x)^2}{\tilde g_t(x)}\,\dd x
=\log\int_{\RR}\tilde g_t(x)\,u_t(x)^2\,\dd x.
\]
又因 \(\int \tilde g_t u_t=\int \tilde h_t=1\)，故
\[
\int \tilde g_t u_t^2
=\int \tilde g_t\bigl((u_t-1)+1\bigr)^2
=1+\int \tilde g_t(u_t-1)^2
=1+\chi^2(\tilde h_t\mid\tilde g_t).
\]
代入命题 \ref{prop:xi-poisson-cauchy-chi2-spectrum-multipole} 即得闭式。
最后使用 Rényi 相对熵关于阶的单调性（\(\alpha\ge1\) 时 \(D_1=D_{\mathrm{KL}}\le D_2\)）得到上界。
\end{proof}

\begin{proposition}[乘极矩的尺度演化链]\label{prop:xi-poisson-cauchy-multipole-moment-chain}\leanverified{paper\_xi\_poisson\_cauchy\_multipole\_moment\_chain}
在命题 \ref{prop:xi-poisson-cauchy-cayley-multipole-fourier} 的设定下，对每个 \(k\ge1\)，函数 \(t\mapsto m_k(t)\) 在 \((0,\infty)\) 上可微，且满足精确微分链
\[
t\,\frac{\dd}{\dd t}m_k(t)=-k\bigl(m_k(t)-m_{k+1}(t)\bigr).
\]
\end{proposition}
\begin{proof}
固定 \(a\in\RR\)，令 \(w(t):=\frac{a}{a+2\mathrm{i}t}\)。
直接求导得
\[
w'(t)=-\frac{2\mathrm{i}a}{(a+2\mathrm{i}t)^2}.
\]
另一方面，
\(
w(t)\,(1-w(t))=\frac{2\mathrm{i}ta}{(a+2\mathrm{i}t)^2}
\)，故 \(t\,w'(t)=-w(t)\,(1-w(t))\)。
于是
\[
\frac{\dd}{\dd t}w(t)^k
=k\,w(t)^{k-1}w'(t)
=-\frac{k}{t}\bigl(w(t)^k-w(t)^{k+1}\bigr).
\]
取 \(a=\Gamma-\bar\gamma\) 并对 \(\Gamma\sim\tilde\nu\) 取期望；由于 \(|w(t)|<1\) 且导数有一致有界支配，可交换微分与期望，从而得到所示微分链。
\end{proof}

\begin{proposition}[\texorpdfstring{$\chi^2$}{chi2} 的精确耗散恒等式]\label{prop:xi-poisson-cauchy-chi2-dissipation-identity}
\leanverified{paper\_xi\_poisson\_cauchy\_chi2\_dissipation\_identity}
令 \(\chi^2(t):=\chi^2(\tilde h_t\mid\tilde g_t)\)。
则成立严格耗散恒等式
\[
\frac{\dd}{\dd t}\bigl(t\,\chi^2(t)\bigr)
=-2\sum_{k\ge1}k\,\bigl|m_k(t)-m_{k+1}(t)\bigr|^2\le 0.
\]
并且等号在某个 \(t>0\) 成立当且仅当 \(\tilde\nu=\delta_{\bar\gamma}\)（从而对所有 \(t>0\) 有 \(\tilde h_t\equiv \tilde g_t\)）。
\end{proposition}
\begin{proof}
由命题 \ref{prop:xi-poisson-cauchy-chi2-spectrum-multipole}，
\(
\chi^2(t)=2\sum_{k\ge1}|m_k(t)|^2
\)，从而
\[
\chi^2{}'(t)=4\sum_{k\ge1}\Re\bigl(m_k'(t)\overline{m_k(t)}\bigr).
\]
代入命题 \ref{prop:xi-poisson-cauchy-multipole-moment-chain} 的 \(m_k'(t)=-(k/t)(m_k-m_{k+1})\)，并整理得到
\[
\frac{\dd}{\dd t}\bigl(t\,\chi^2(t)\bigr)
=-2\sum_{k\ge1}k\,|m_k-m_{k+1}|^2.
\]
若右端在某个 \(t>0\) 为零，则 \(m_k(t)=m_{k+1}(t)\) 对所有 \(k\ge1\) 成立。
又由 \(|m_k(t)|\le \EE|W_t(\Gamma)|^k\to 0\)（\(k\to\infty\)）可得 \(m_k(t)=0\) 对所有 \(k\ge1\)。
由命题 \ref{prop:xi-poisson-cauchy-cayley-multipole-fourier} 的傅里叶展开得到 \(u_t\equiv 1\)，即 \(\tilde h_t\equiv \tilde g_t\)，从而 \(\tilde\nu=\delta_{\bar\gamma}\)。
反向蕴含显然。
\end{proof}

\begin{proposition}[尺度 Möbius 共变与 Koenigs 线性化]\label{prop:xi-poisson-cauchy-multipole-mobius-koenigs}
\leanverified{paper\_xi\_poisson\_cauchy\_multipole\_mobius\_koenigs}
对任意 \(s>0\) 定义 Möbius 自映射
\[
\Phi_s(u):=\frac{u}{s-(s-1)u}\qquad(u\in\mathbb D).
\]
则对每个 \(\gamma\in\RR\) 与 \(t>0\) 有严格共变关系
\[
W_{st}(\gamma)=\Phi_s\bigl(W_t(\gamma)\bigr),
\]
并且 \((\Phi_s)_{s>0}\) 构成乘法参数化半群：\(\Phi_s\circ\Phi_r=\Phi_{sr}\)。
再者，Koenigs 坐标
\[
\kappa(u):=\frac{u}{1-u}
\]
将该半群对角化为纯伸缩：
\[
\kappa\!\left(W_t(\gamma)\right)=\frac{\gamma-\bar\gamma}{2\mathrm{i}t},
\qquad
\kappa(\Phi_s(u))=\frac{1}{s}\,\kappa(u).
\]
\end{proposition}
\begin{proof}
令 \(a:=\gamma-\bar\gamma\)，则 \(W_t(\gamma)=\frac{a}{a+2\mathrm{i}t}\)。
由此解出 \(a=2\mathrm{i}t\,\frac{W_t}{1-W_t}\)，从而
\[
W_{st}
=\frac{a}{a+2\mathrm{i}st}
=\frac{\frac{W_t}{1-W_t}}{\frac{W_t}{1-W_t}+s}
=\frac{W_t}{s-(s-1)W_t}
=\Phi_s(W_t).
\]
半群性质由直接代入或由 Koenigs 坐标下的线性化得到。
最后由 \(\kappa(W_t)=\frac{W_t}{1-W_t}=\frac{a}{2\mathrm{i}t}\) 得到对角化恒等式。
\end{proof}

\begin{proposition}[有限模态压缩的 \texorpdfstring{$L^2(\tilde g_t)$}{L2(gt)} 精确误差]\label{prop:xi-poisson-cauchy-multipole-truncation-l2}
\leanverified{paper\_xi\_poisson\_cauchy\_multipole\_truncation\_l2}
对 \(N\ge1\) 定义截断
\[
u_{t,N}(\zeta):=1+\sum_{k=1}^{N}\Bigl(m_k(t)\,\zeta^{-k}+\overline{m_k(t)}\,\zeta^{k}\Bigr),
\qquad \zeta\in\TT,
\]
并令 \( \tilde h_{t,N}(x):=\tilde g_t(x)\,u_{t,N}\bigl(\zeta_t(x)\bigr)\)。
则有严格等式
\[
\int_{\RR}\frac{\bigl(\tilde h_t(x)-\tilde h_{t,N}(x)\bigr)^2}{\tilde g_t(x)}\,\dd x
=2\sum_{k>N}|m_k(t)|^2.
\]
从而 \((m_1(t),\dots,m_N(t))\) 构成 \(L^2(\tilde g_t)\) 意义下最优 \(N\)-模态近似的完备信息，残差能量完全由高频乘极能量尾项给出。
\end{proposition}
\begin{proof}
在 \(\TT\) 上以 Haar 内积，\(u_{t,N}\) 为 \(u_t\) 在次数不超过 \(N\) 的三角多项式子空间上的正交投影；
故
\[
\|u_t-u_{t,N}\|_{L^2(\TT)}^2=\sum_{|k|>N}|c_k|^2=2\sum_{k>N}|m_k(t)|^2.
\]
再由 \(\tilde g_t(x)\,\dd x\) 为 Haar 测度回拉得
\(
\|u_t-u_{t,N}\|_{L^2(\TT)}^2=\int(\tilde h_t-\tilde h_{t,N})^2/\tilde g_t
\)。
\end{proof}

\endinput


\begin{conclusion}[\texorpdfstring{$t^{-6}$}{t-6} 系数的严格负性与 KL 主系数的下逼近]\label{con:xi-poisson-kl-sixth-coefficient-strict-negativity}
在结论 \ref{con:xi-poisson-kl-two-term-sharp-explicit} 的矩条件下，令
\[
B_6:=\frac{\sigma^6}{64}-\frac{\sigma^2\mu_4}{8}+\frac{3\mu_3^2}{32}.
\]
则 \(B_6<0\)。从而当 \(t\to\infty\) 时有
\[
t^4D_{\mathrm{KL}}(\tilde h_t\mid \tilde g_t)=\frac{\sigma^4}{8}+\frac{B_6}{t^2}+o(t^{-2}),
\]
并且在充分大 \(t\) 上严格小于 \(\sigma^4/8\)。
进一步地，函数 \(t\mapsto t^4D_{\mathrm{KL}}(\tilde h_t\mid \tilde g_t)\) 最终严格单调递增并收敛到 \(\sigma^4/8\)。
\end{conclusion}
\begin{proof}[证明]
由中心矩不等式
\[
\mu_4\ge \sigma^4,\qquad
\mu_3^2\le \sigma^2\mu_4
\]
可得
\[
B_6
\le \frac{\sigma^6}{64}-\frac{\sigma^2\mu_4}{8}+\frac{3\sigma^2\mu_4}{32}
=\frac{\sigma^6}{64}-\frac{\sigma^2\mu_4}{32}
\le \frac{\sigma^6}{64}-\frac{\sigma^6}{32}
=-\frac{\sigma^6}{64}<0.
\]
若两次不等式同时取等，则需同时满足 \(\mu_4=\sigma^4\) 与 \(\mu_3^2=\sigma^2\mu_4\)，但前者迫使 \(|\gamma-\bar\gamma|\) 几乎处处常值，从而 \(\mu_3=0\)，与后者矛盾，故严格性成立。

由结论 \ref{con:xi-poisson-kl-two-term-sharp-explicit} 的展开立得
\(
t^4D_{\mathrm{KL}}=\sigma^4/8+B_6t^{-2}+o(t^{-2})
\)，因此 \(t\) 充分大时严格小于 \(\sigma^4/8\)。
对该展开求导得到
\[
\frac{\dd}{\dd t}\Bigl(t^4D_{\mathrm{KL}}(\tilde h_t\mid \tilde g_t)\Bigr)
=-\frac{2B_6}{t^3}+o(t^{-3})>0
\]
在充分大 \(t\) 上成立，从而得到最终严格单调递增及收敛。证毕。
\end{proof}

\begin{conclusion}[光滑 \texorpdfstring{$f$}{f}-散度的两项锐渐近]\label{con:xi-poisson-fdivergence-two-term-sharp}
在结论 \ref{con:xi-poisson-kl-two-term-sharp-explicit} 的矩条件下，令 \(f:(0,\infty)\to\RR\) 为凸函数，满足 \(f(1)=0\)，并在 \(1\) 的邻域三阶可微且 \(f''(1),f'''(1)\) 有限。
定义
\[
D_f(\tilde h_t\mid \tilde g_t):=\int_{\RR}\tilde g_t(x)\,f\!\left(\frac{\tilde h_t(x)}{\tilde g_t(x)}\right)\dd x.
\]
则当 \(t\to\infty\) 有
\[
D_f(\tilde h_t\mid \tilde g_t)
=\frac{f''(1)}{8}\frac{\sigma^4}{t^4}
+\frac1{t^6}\left[
f''(1)\!\left(\frac{3\mu_3^2}{32}-\frac{\sigma^2\mu_4}{8}\right)
-\frac{f'''(1)}{64}\sigma^6
\right]
+o(t^{-6}).
\]
特别地，至 \(t^{-6}\) 阶为止不出现奇次幂项。
\end{conclusion}
\begin{proof}[证明]
令 \(\delta_t:=\tilde h_t/\tilde g_t-1\)。由四阶正规形与尺度律（见定理 \ref{thm:conclusion-poisson-cauchy-kl-sixth-order-expansion} 的证明结构）有 \(\delta_t=O(t^{-2})\) 且 \(\int_{\RR}\tilde g_t\,\delta_t\,\dd x=0\)。
对 \(f(1+\delta)\) 在 \(\delta=0\) 处作三阶 Taylor 展开并在 \(\tilde g_t\dd x\) 下积分，利用一次项消失得到
\[
D_f(\tilde h_t\mid \tilde g_t)
=\frac{f''(1)}{2}\int_{\RR}\tilde g_t\,\delta_t^2\,\dd x
+\frac{f'''(1)}{6}\int_{\RR}\tilde g_t\,\delta_t^3\,\dd x
+o(t^{-6}).
\]
将 \(\delta_t\) 的 \(t^{-4}\) 精细展开代回并按阶收集；由于二阶轮廓为偶、三阶轮廓为奇，导致 \(t^{-5}\) 的混合项积分严格为零，从而仅保留偶次幂至此阶。
所需有理积分常数与 KL 计算一致，代回即得所示系数。证毕。
\end{proof}

\begin{conclusion}[耗散量 \texorpdfstring{$I_1$}{I1} 的七阶锐渐近]\label{con:xi-poisson-i1-two-term-sharp}
在结论 \ref{con:xi-poisson-kl-two-term-sharp-explicit} 的假设下，若 \(I_1(\tilde h_t\mid \tilde g_t)\) 满足耗散恒等式
\[
\frac{\dd}{\dd t}D_{\mathrm{KL}}(\tilde h_t\mid \tilde g_t)=-I_1(\tilde h_t\mid \tilde g_t),
\]
则当 \(t\to\infty\) 有
\[
I_1(\tilde h_t\mid \tilde g_t)
=\frac{\sigma^4}{2t^5}
+\frac{6B_6}{t^7}
+o(t^{-7}),
\]
其中 \(B_6\) 如结论 \ref{con:xi-poisson-kl-sixth-coefficient-strict-negativity} 所定义。
此外，当 \(\sigma^2>0\) 时，\(t^5 I_1(\tilde h_t\mid \tilde g_t)\) 最终严格单调递增并收敛到 \(\sigma^4/2\)。
\end{conclusion}
\begin{proof}[证明]
将结论 \ref{con:xi-poisson-kl-two-term-sharp-explicit} 的展开逐项求导并取负号即可得到 \(I_1\) 的两项展开。
由结论 \ref{con:xi-poisson-kl-sixth-coefficient-strict-negativity} 知 \(B_6<0\)，从而
\[
t^5I_1(\tilde h_t\mid \tilde g_t)=\frac{\sigma^4}{2}+\frac{6B_6}{t^2}+o(t^{-2})
\]
从下方逼近 \(\sigma^4/2\)，并同理由导数符号推出最终严格单调递增。证毕。
\end{proof}

\begin{conclusion}[六阶塌缩：\texorpdfstring{$f$}{f}-散度的测度依赖由单一标量不变量控制]\label{con:xi-poisson-fdivergence-single-parameter-collapse}
在结论 \ref{con:xi-poisson-fdivergence-two-term-sharp} 的条件下，定义
\[
\mathfrak J(\tilde\nu):=4\sigma^2\mu_4-3\mu_3^2.
\]
则当 \(\sigma^2>0\) 时有 \(\mathfrak J(\tilde\nu)>0\)，并且对任意满足结论 \ref{con:xi-poisson-fdivergence-two-term-sharp} 条件的 \(f\) 有
\[
D_f(\tilde h_t\mid \tilde g_t)
=\frac{f''(1)}{8}\frac{\sigma^4}{t^4}
-\frac{f''(1)}{32}\frac{\mathfrak J(\tilde\nu)}{t^6}
-\frac{f'''(1)}{64}\frac{\sigma^6}{t^6}
+o(t^{-6}),
\qquad t\to\infty.
\]
特别地，至 \(t^{-6}\) 阶为止，\(f\)-散度对 \(\tilde\nu\) 的非高斯信息仅通过 \(\mathfrak J(\tilde\nu)\) 发生依赖；若 \(\tilde\nu\) 关于 \(\bar\gamma\) 对称（从而 \(\mu_3=0\)），则 \(\mathfrak J(\tilde\nu)=4\sigma^2\mu_4\)，从而在 \(t^{-4}\) 项已确定 \(\sigma^2\) 后，任取一类 \(f\)-散度的 \(t^{-6}\) 系数即可唯一确定 \(\mu_4\)。
\end{conclusion}
\begin{proof}[证明]
由结论 \ref{con:xi-poisson-fdivergence-two-term-sharp}，
\[
\frac{3\mu_3^2}{32}-\frac{\sigma^2\mu_4}{8}
=-\frac1{32}\bigl(4\sigma^2\mu_4-3\mu_3^2\bigr)
=-\frac1{32}\mathfrak J(\tilde\nu),
\]
代回即可得到塌缩表达式。
又由 \(\mu_3^2\le \sigma^2\mu_4\) 与 \(\mu_4>0\) 可知当 \(\sigma^2>0\) 时 \(\mathfrak J(\tilde\nu)>0\)。证毕。
\end{proof}

\begin{conclusion}[Poisson--Cauchy 粗粒化误差的规范化 \texorpdfstring{$L^p$}{Lp} 锐极限]\label{con:xi-poisson-cauchy-lp-sharp-limit-canonical-profile}
在结论 \ref{con:xi-poisson-lp-spectrum-sharp-family} 的条件下，记
\[
\mu_2:=\sigma^2,\qquad
R(y):=\frac1\pi\frac{3y^2-1}{(1+y^2)^3}.
\]
则对任意 \(p\in(1,\infty]\) 有
\[
\lim_{t\to\infty}
\frac{t^{3-\frac1p}}{\mu_2}\,
\|\tilde h_t-\tilde g_t\|_{L^p(\RR)}
=\|R\|_{L^p(\RR)}.
\]
并且
\[
\|R\|_{L^2(\RR)}=\frac{\sqrt3}{4\sqrt\pi},\qquad
\|R\|_{L^\infty(\RR)}=\frac1\pi.
\]
\end{conclusion}
\begin{proof}
由结论 \ref{con:xi-poisson-lp-spectrum-sharp-family}，
\[
\lim_{t\to\infty}
\frac{t^{3-\frac1p}}{\sigma^2}\,
\|\tilde h_t-\tilde g_t\|_{L^p}
=\frac1\pi\left(\int_{\RR}\frac{|3y^2-1|^p}{(1+y^2)^{3p}}\,\dd y\right)^{1/p}
\]
（\(p<\infty\)），且 \(p=\infty\) 端点极限为 \(1/\pi\)。
右端即 \(\|R\|_{L^p}\) 的定义形式，故第一式成立。
对 \(p=2\)，
\[
\|R\|_{L^2}^2
=\frac1{\pi^2}\int_{\RR}\frac{(3y^2-1)^2}{(1+y^2)^6}\,\dd y
=\frac{3}{16\pi},
\]
从而 \(\|R\|_{L^2}=\frac{\sqrt3}{4\sqrt\pi}\)。
又 \(\|R\|_\infty=1/\pi\) 已由同一结论给出。证毕。
\end{proof}

\begin{conclusion}[相对熵二阶校正的 \texorpdfstring{$\mu_2$}{mu2}-重参数化显式系数]\label{con:xi-poisson-kl-second-correction-mu2-notation}
在结论 \ref{con:xi-poisson-kl-two-term-sharp-explicit} 的矩条件下，记
\[
\mu_k:=\int_{\RR}(\gamma-\bar\gamma)^k\,\dd\tilde\nu(\gamma)\qquad(k=2,3,4),\qquad \mu_2=\sigma^2.
\]
则
\[
D_{\mathrm{KL}}(\tilde h_t\mid \tilde g_t)
=\frac{\mu_2^2}{8t^4}
+\frac1{t^6}\left(\frac{3\mu_3^2}{32}-\frac{\mu_2\mu_4}{8}+\frac{\mu_2^3}{64}\right)
+o(t^{-6}).
\]
\end{conclusion}
\begin{proof}
将结论 \ref{con:xi-poisson-kl-two-term-sharp-explicit} 中的 \(\sigma^2\) 统一替换为 \(\mu_2\) 即得。
其 \(t^{-6}\) 系数来自 \(\delta_t\)-展开中的二次与三次项积分常数
\[
\int Gu_2^2=\frac14,\quad
\int Gu_3^2=\frac{3}{16},\quad
\int Gu_2u_4=-\frac18,\quad
\int Gu_2^3=-\frac{3}{32},
\]
与该结论的推导完全一致。证毕。
\end{proof}

\begin{conclusion}[任意光滑 \texorpdfstring{$f$}{f}-散度的二阶校正与 \texorpdfstring{$f'''(1)$}{f'''(1)} 非普适项]\label{con:xi-poisson-fdiv-second-correction-f3-term}
在结论 \ref{con:xi-poisson-fdivergence-two-term-sharp} 的假设下，
\[
D_f(\tilde h_t\mid\tilde g_t)
=\frac{f''(1)}{8}\frac{\mu_2^2}{t^4}
+\frac1{t^6}\left[
f''(1)\!\left(\frac{3\mu_3^2}{32}-\frac{\mu_2\mu_4}{8}\right)
-\frac{f'''(1)}{64}\mu_2^3
\right]
+o(t^{-6}).
\]
\end{conclusion}
\begin{proof}
由结论 \ref{con:xi-poisson-fdivergence-two-term-sharp} 直接作 \(\mu_2=\sigma^2\) 重参数化即可得到所示式。
其中 \(f''(1)\)-项由 \(\int \tilde g_t\delta_t^2\) 的 \(t^{-6}\) 系数贡献，
\(f'''(1)\)-项由 \(\int \tilde g_t\delta_t^3\) 的首项贡献，故该项不具有普适性。证毕。
\end{proof}

\begin{conclusion}[Cayley--Poisson 融合恒等式与双侧核界]\label{con:xi-cayley-poisson-scale-integral-identity-bounds}
令 Cayley 变换 \(\mathcal C(z):=(z-\mathrm{i})/(z+\mathrm{i})\)，并取
\[
w=\mathcal C(\gamma+\mathrm{i}\delta),\qquad \gamma\in\RR,\ \delta>0.
\]
则
\[
|w|^2=\frac{\gamma^2+(\delta-1)^2}{\gamma^2+(\delta+1)^2},
\]
并有尺度积分恒等式
\[
-\log|w|
=\int_{|1-\delta|}^{1+\delta}\frac{a}{\gamma^2+a^2}\,\dd a
=\pi\int_{|1-\delta|}^{1+\delta}\mathsf P_a(\gamma)\,\dd a,
\qquad
\mathsf P_a(\gamma):=\frac1\pi\frac{a}{\gamma^2+a^2}.
\]
进一步，成立双侧界
\[
\frac{2\delta}{\gamma^2+(1+\delta)^2}
\le
-\log|w|
\le
\frac{2\delta}{\gamma^2+(|1-\delta|)^2}.
\]
\end{conclusion}
\begin{proof}
由 \(|w|^2\) 公式，
\[
-2\log|w|
=\log\!\bigl(\gamma^2+(\delta+1)^2\bigr)-\log\!\bigl(\gamma^2+(\delta-1)^2\bigr).
\]
设 \(F(a):=\log(\gamma^2+a^2)\)，则 \(F'(a)=2a/(\gamma^2+a^2)\)，故
\[
-2\log|w|
=\int_{\delta-1}^{\delta+1}\frac{2a}{\gamma^2+a^2}\,\dd a
=\int_{|1-\delta|}^{1+\delta}\frac{2a}{\gamma^2+a^2}\,\dd a,
\]
从而得到尺度积分恒等式。
再写作
\[
-2\log|w|
=\log\!\left(1+\frac{4\delta}{\gamma^2+(\delta-1)^2}\right),
\]
并对 \(u>0\) 使用
\(
\frac{u}{1+u}\le\log(1+u)\le u
\)，即得双侧界。证毕。
\end{proof}

\begin{conclusion}[Jensen 缺陷的 Poisson 足迹下界]\label{con:xi-jensen-defect-poisson-footprint-lower-bound}
沿用 \eqref{eq:xi-comoving-shifted-Fzeta}--\eqref{eq:xi-comoving-shifted-jensen-defect} 的记号。
若 \(\rho=\beta+\mathrm{i}\gamma\) 为离临界线零点并记 \(\delta:=\frac12-\beta>0\)，令
\[
a(\rho;x_0):=\mathcal C\bigl((\gamma-x_0)+\mathrm{i}\delta\bigr),
\qquad
|a|^2=\frac{(\gamma-x_0)^2+(\delta-1)^2}{(\gamma-x_0)^2+(\delta+1)^2}.
\]
则对任意 \(x_0\in\RR\) 有
\[
\liminf_{\varrho\to1^-}D_{\zeta,x_0}(\varrho)
\ge
\log\frac{1}{|a(\rho;x_0)|}
\ge
\frac{2\delta}{(\gamma-x_0)^2+(1+\delta)^2}
=2\pi\frac{\delta}{1+\delta}\,\mathsf P_{1+\delta}(\gamma-x_0).
\]
\end{conclusion}
\begin{proof}
由命题 \ref{prop:app-jensen-single-zero-lower-bound}，对 \(F_{\zeta,x_0}\) 的任一盘内零点 \(a\) 有
\[
\liminf_{\varrho\to1^-}D_{\zeta,x_0}(\varrho)\ge \log\frac1{|a|}.
\]
将 \(a=a(\rho;x_0)\) 代入并应用结论
\ref{con:xi-cayley-poisson-scale-integral-identity-bounds}
的下界（取 \(\gamma\mapsto \gamma-x_0\)）得到
\[
\log\frac1{|a(\rho;x_0)|}
\ge \frac{2\delta}{(\gamma-x_0)^2+(1+\delta)^2}.
\]
最后由
\(
\mathsf P_{1+\delta}(u)=\frac1\pi\frac{1+\delta}{u^2+(1+\delta)^2}
\)
直接化简得到 Poisson 足迹形式。证毕。
\end{proof}

